\begin{exercise}[Why charged vectors must be gauge fields]{I-CH08-002}
Assume a local, Lorentz-invariant, unitary four-dimensional theory with a
weakly coupled perturbative expansion and finitely many spin-one particles.
Let a vector $W_\mu$ carry a nontrivial charge under a compact gauge group.
Analyze the ultraviolet behavior of its propagator and of amplitudes with
longitudinal external polarizations.  Show that cancellation of the terms
which grow as powers of $E/m_W$ fixes the cubic and quartic vector couplings
to Yang--Mills form.  Derive the antisymmetry and Jacobi conditions on the
cubic couplings, and explain how a mass for $W_\mu$ can then arise through a
Higgs or Stueckelberg description.  Conclude, under the stated assumptions,
that a perturbatively renormalizable charged vector is a gauge connection of
an enlarged gauge group rather than an independent vector matter field.
\end{exercise}
